测得的特性曲线表现为[填写：线性/部分非线性]行为。回归直线为 $U = [m] \cdot T + [b]$。相关系数 $r = [\,]$ 表明[评价]。参考温度与基于回归计算的温度之间的最大偏差为 $[\,]\,\si{\celsius}$，出现在 $[\,]\,\si{\celsius}$ 处。

绝对偏差的平均值为 $[\,]\,\si{\celsius}$。观察到的离散程度为 $[\,]\,\si{\celsius}$。因此，实际精度比仅由 ADC 量化所预期的分辨率[更好/更差]。

由测得的幅值得到 $A_{\mathrm{dB}} = 20 \log_{10}(U_{\mathrm{nach}}/U_{\mathrm{vor}}) = [\,]\,\si{\decibel}$。因此要求的 $[\,]\,\si{\decibel}$ 衰减[达到/未达到]。数字温度显示的剩余波动为 $[\,]\,\si{\celsius}$。由于时间常数 $\tau = [\,]\,\si{\second}$，产生了可观察到的延迟：[观察结果]。

附加放大显著改善了 ADC 输入范围的利用。在理想情况下，温度等效量化步长从约 $\SI{1.96}{\celsius}$ 降至约 $\SI{0.47}{\celsius}$。实验确定的斜率 $[\,]\,\si{\volt\per\celsius}$ 与理论值 $[\,]\,\si{\volt\per\celsius}$ 相差 $[\,]\,\%$。可能原因包括电阻公差、AD595 的实际增益、偏置电压以及电压测量精度有限。

校准特性曲线应根据 $r$、$R^2$ 和残差进行评价。较高的相关系数虽然确认了线性趋势，但不能说明某些温度区间是否存在系统偏差。因此尤其需要考察差值 $T_{\mathrm{Anzeige}} - T_{\mathrm{ref}}$。在本实验中，最大绝对偏差为 $[\,]\,\si{\celsius}$。[此处评价该偏差对于任务是否可接受。]

RC 低通滤波器将高频干扰降低了 $[\,]\,\si{\decibel}$。理论衰减与实测衰减之间的差值为 $[\,]\,\si{\decibel}$。可能原因包括 $R$ 和 $C$ 的公差、前级有限的输出阻抗、测量卡输入负载以及干扰并非纯正弦信号。同时，低通滤波器会减慢温度变化过程。因此所选截止频率是在抑制干扰和响应速度之间的折中。
